% id: 138-01
% book: 베이직쎈 미적분1 · 09 정적분의 활용
% section: 개념 56 곡선과 $x$축 사이의 넓이
% passage: 다음 그림에서 색칠한 부분의 넓이를 구하시오.
% figure: tikz:fig-138-01
% answer: 
% answer_source: 
% variant_level: 0 (원본 전사)
% 이 파일은 build.mjs 가 items.json 에서 만든다. 고칠 때는 items.json 을 고친다.
\smash{\raisebox{1.5mm}{\dmkichul{베이직쎈 미적분1 138쪽 01번}}}
\noindent\begin{minipage}[t]{0.58\linewidth}\vspace{0pt}\raggedright \end{minipage}\hfill\begin{minipage}[t]{0.4\linewidth}\vspace{0pt}\centering \resizebox{\ifdim\width>\linewidth\linewidth\else\width\fi}{!}{% 138-01(138쪽) — 포물선 $y=x^2-2x+2$ 와 $x$축·직선 $x=2$ 로 둘러싸인 색칠 부분. 스캔 화질이 낮아 TikZ 로 다시 그림.
% 원문에 있는 것만 옮긴다: 곡선 · 색칠 부분 · 오른쪽 경계 세로선 · 라벨 $y=x^2-2x+2$, 2($y$축), 2($x$축), O, x, y. 원문에 없는 눈금·넓이 값은 적지 않는다.
\begin{tikzpicture}[x=0.6cm, y=0.46cm, line width=0.6pt, font=\small]
  \fill[gray!20] (0,0) -- plot[domain=0:2, samples=60] (\x, {(\x)*(\x)-2*(\x)+2}) -- (2,0) -- cycle;
  \draw[->, >=stealth, line width=0.7pt] (-0.8,0) -- (3.2,0) node[below=1pt] {$x$};
  \draw[->, >=stealth, line width=0.7pt] (0,-0.95) -- (0,3.5) node[left=1pt] {$y$};
  \draw[domain=-0.3:2.3, samples=90, smooth] plot (\x, {(\x)*(\x)-2*(\x)+2});
  \draw (2,0) -- (2,2);
  \node[below left=-2pt and -1pt] at (0,0) {$\pt{O}$};
  \node[left=2pt] at (0,2) {$2$};
  \node[below=2pt] at (2,0) {$2$};
  \node[anchor=west, inner sep=1pt] at (0.45,3.25) {$y=x^2-2x+2$};
\end{tikzpicture}
}\end{minipage}\par
\dmhint{$\displaystyle\int_{0}^{\text{\scriptsize\setlength{\fboxsep}{1.5pt}\framebox[1.5em]{\rule{0pt}{1.6ex}}}}|x^2-2x+2|\mkern3mu dx$\\ $\displaystyle =\int_{0}^{\text{\scriptsize\setlength{\fboxsep}{1.5pt}\framebox[1.5em]{\rule{0pt}{1.6ex}}}}(x^2-2x+2)\mkern3mu dx$\\ $=\Big[\dfrac{1}{3}x^3-x^2+2x\Big]_{0}^{\text{\scriptsize\setlength{\fboxsep}{1.5pt}\framebox[1.5em]{\rule{0pt}{1.6ex}}}}$\\ $=\blank$}
